\begin{figure}[H]
\centering
\resizebox{\textwidth}{!}{%
\begin{tikzpicture}[
    node distance=1.2cm and 2.5cm,
    >=Stealth,
    thick,
    font=\sffamily\large
]

% ---------------------------------------------------------
% Styles
% ---------------------------------------------------------
\tikzset{
    op_node/.style={
        circle, draw=black, very thick, minimum size=0.9cm, fill=white
    },
    io_node/.style={
        font=\bfseries\large
    },
    dot/.style={
        circle, fill=black, inner sep=0pt, minimum size=4pt
    }
}

% ---------------------------------------------------------
% 1. Multipliers (The Column Backbone)
% ---------------------------------------------------------
\node[op_node] (Mult1) at (0, 4.5) {$\times$};
\node[op_node] (Mult2) at (0, 1.5) {$\times$};
\node[op_node] (Mult3) at (0, -1.5) {$\times$};
\node[op_node] (Mult4) at (0, -4.5) {$\times$};

% Labels
\node[above right=0.1cm of Mult1, font=\scriptsize] {$Re \cdot \cos$};
\node[above right=0.1cm of Mult2, font=\scriptsize] {$Im \cdot \sin$};
\node[above right=0.1cm of Mult3, font=\scriptsize] {$Re \cdot \sin$};
\node[above right=0.1cm of Mult4, font=\scriptsize] {$Im \cdot \cos$};

% ---------------------------------------------------------
% 2. Inputs (Far Left)
% ---------------------------------------------------------
\node[io_node] (InRe)  at (-5, 3.0) {$x_{Re}$};
\node[io_node] (InCos) at (-5, 1.0) {$\cos \theta$};
\node[io_node] (InSin) at (-5, -1.0) {$\sin \theta$};
\node[io_node] (InIm)  at (-5, -3.0) {$x_{Im}$};

% ---------------------------------------------------------
% 3. Orthogonal Wiring
% ---------------------------------------------------------
% Re -> Mult1 and Mult3
\draw (InRe) -- (-3.5, 3.0) coordinate (SplitRe);
\node[dot] at (SplitRe) {};
\draw[->] (SplitRe) |- (Mult1.west);
\draw[->] (SplitRe) |- (Mult3.west);

% Im -> Mult2 and Mult4
\draw (InIm) -- (-3.0, -3.0) coordinate (SplitIm);
\node[dot] at (SplitIm) {};
\draw[->] (SplitIm) |- (Mult2.west);
\draw[->] (SplitIm) |- (Mult4.west);

% Cos -> Mult1 and Mult4
\draw (InCos) -- (-2.5, 1.0) coordinate (SplitCos);
\node[dot] at (SplitCos) {};
\draw[->] (SplitCos) |- (Mult1.150);
\draw[->] (SplitCos) |- (Mult4.210);

% Sin -> Mult2 and Mult3
\draw (InSin) -- (-2.0, -1.0) coordinate (SplitSin);
\node[dot] at (SplitSin) {};
\draw[->] (SplitSin) |- (Mult2.210);
\draw[->] (SplitSin) |- (Mult3.150);

% ---------------------------------------------------------
% 4. Output Stage
% ---------------------------------------------------------
% Subtractor
\node[op_node] (Sub) at (3, 3.0) {$-$};
% Adder
\node[op_node] (Add) at (3, -3.0) {$+$};

\draw[->] (Mult1) -| (Sub.north);
\draw[->] (Mult2) -| (Sub.south);
\draw[->] (Mult3) -| (Add.north);
\draw[->] (Mult4) -| (Add.south);

% ---------------------------------------------------------
% 5. Final Outputs
% ---------------------------------------------------------
\node[io_node, right=1.0cm of Sub] (OutI) {$I_{out}$};
\node[io_node, right=1.0cm of Add] (OutQ) {$Q_{out}$};

\draw[->] (Sub) -- (OutI);
\draw[->] (Add) -- (OutQ);

\end{tikzpicture}%
}

\vspace{5mm}
\begin{align*}
    I_{out} &= (x_{Re} \cdot \cos\theta) - (x_{Im} \cdot \sin\theta) \\
    Q_{out} &= (x_{Re} \cdot \sin\theta) + (x_{Im} \cdot \cos\theta)
\end{align*}

\noindent\textit{All four products computed in parallel (64-bit). Result truncated to $[47\!:\!16]$ (Q16.16). Enable gate provides sample-and-hold behavior.}

\caption{Complex\_Mult\_Output block -- complex multiplication with orthogonal signal routing.}
\label{fig:complex_mult}
\end{figure}
